%-----------------------------------LICENSE------------------------------------%
%   This file is free software: you can redistribute it and/or                 %
%   modify it under the terms of the GNU General Public License as             %
%   published by the Free Software Foundation, either version 3 of the         %
%   License, or (at your option) any later version.                            %
%                                                                              %
%   This file is distributed in the hope that it will be useful,               %
%   but WITHOUT ANY WARRANTY; without even the implied warranty of             %
%   MERCHANTABILITY or FITNESS FOR A PARTICULAR PURPOSE.  See the              %
%   GNU General Public License for more details.                               %
%                                                                              %
%   You should have received a copy of the GNU General Public License along    %
%   with this file.  If not, see <https://www.gnu.org/licenses/>.              %
%------------------------------------------------------------------------------%
%       Author: Ryan Maguire                                                   %
%       Date:   2023/10/14                                                     %
%------------------------------------------------------------------------------%
\documentclass{article}
\usepackage[margin=1in]{geometry}
\usepackage{amssymb}
\usepackage{amsthm}
\usepackage{amsmath}
\usepackage{hyperref}
\hypersetup{colorlinks=true, linkcolor=blue}
\title{Research Statement}
\author{Ryan Maguire}
\date{\today}

\begin{document}
    \maketitle
    \section{Introduction}
        My research interests lie in astronomy, numerical mathematics,
        and spacetimes.
        I specialize in solving complex inverse problems by developing
        highly optimized numerical methods and parallelized software.
        My work has allowed researchers to analyze the ringed planets in our
        solar system to much higher resolutions than previously achievable,
        and my goal is to extend this work to the stars beyond.
    \section{Planetary Rings}
        Over the last 9 years I have been working with Richard French from
        Wellesley college analyzing the radio science data from the NASA
        Cassini mission and the Voyager spacecraft. Radio waves from these
        spacecrafts pass through the rings of Saturn and Uranus on their way
        to Earth creating \textit{occultation observations}
        \cite{Marouf1982TheoryOR}. These radio waves diffract as they pass
        through the rings, and this diffracted data is measured on Earth. We
        may model this using the Rayleigh-Sommerfeld equation:
        \begin{equation}
            \hat{T}(\rho_{0},\,\phi_{0})
            =
            \frac{\sin(B)}{i\lambda}
            \int_{0}^{\infty}
                \int_{0}^{2\pi}
                    \frac{\rho}{D}
                    T(\rho,\,\phi)
                    \exp\big(i\psi(\rho,\,\phi,\,\rho_{0},\,\phi_{0})\big)\;
                \textrm{d}\rho\;
                \textrm{d}\phi
        \end{equation}
        where $\lambda$ is the wavelength of the incident light, $B$ is
        the angle the ray makes with the ring plane, $(\rho_{0},\,\phi_{0})$
        is the point of interest, $(\rho,\,\phi)$ is the dummy point of
        integration, $T$ is the transmittance of the rings, $\psi$ is
        the Fresnel phase, and $D$ is the distance from the spacecraft to
        the point $(\rho_{0},\,\phi_{0})$. Following \cite{MTR86}, we reduce
        this down to a single integral using the stationary phase approximation,
        and then obtain an approximate inverse formula by swapping the variables
        of integration and taking the complex conjugate of the integrand.
        If $\phi_{s}$ denotes the stationary azimuth angle
        (the angle with $\partial\psi/\partial\phi=0$), and if $\psi_{s}$
        denotes the evaluation of $\psi$ at the tuple
        $(\rho,\,\phi_{s},\,\rho_{0},\,\phi_{0})$, then this approximate inverse
        may be written:
        \begin{equation}
            T(\rho)
            \approx
            \frac{\sin(B)}{i\lambda}
            \int_{0}^{\infty}
                \hat{T}(\rho_{0})
                \sqrt{\frac{2\pi}{|\psi_{s}^{\prime\prime}|}}
                \frac{\exp\left(-i(\psi_{s}-\pi/4)\right)}{D}\,
                    \textrm{d}\rho_{0}
        \end{equation}
        We implement this in \texttt{rss\_ringoccs}
        \cite{rssringoccs}, a suite of tools in C and
        Python to numerically solve this inverse problem,
        obtaining the function $T(\rho,\,\phi)$ from an observed
        $\hat{T}(\rho_{0},\,\phi_{0})$ profile.
        \par\hfill\par
        Over the years we have explored
        FFT methods, new quadrature techniques, forward modeling of diffraction
        using ideal profiles, and interpolation methods to accurately and
        efficiently work with the geometric data in an occultation observation.
        This has been highly optimized to very efficiently process the data.
        All of the occultation observations from Cassini
        can be processed at a resolution of 1000 meters in about 20 minutes on
        a 2017 8-core iMac. The initial python code from 2017 took 130+
        hours on that same computer, but a large effort was undergone between
        2019 and 2020 to reimplement everything in C.
        In early 2026 I underwent the effort of ensuring
        thread-safety and reentrancy for the routines implemented in
        \cite{libtmpl} and \cite{rssringoccs} to allow safe parallel processing.
        This effort was successful, and a 32-core machine can process an entire
        occultation observation accurately in a few seconds.
        Benchmarks for processing the entirety of the Cassini data will be
        available soon.
        \par\hfill\par
        All of this has been used to study the rings and led to
        \cite{FRENCH2023115678} and \cite{NICHOLSON2023115287}.
        During the summer of 2025 I developed new numerical techniques to
        process the diffracted data very carefully, allowing us to obtain new
        ring profiles at incredibly high resolutions. We've been able to push
        the methods all the way down to 50 or 20 or even 10 meters, depending on the
        data set, at which point the noise in the measurements start to become
        an issue. We are now seeing new
        structures in the rings of Saturn for the first time. We are also now
        actively returning to the Voyager data for Saturn and Uranus to
        see if these new techniques can reveal anything new. The numerical
        methods and the results obtained are currently being compiled into a
        paper, and we hope to submit these efforts for publication
        in the near future.
    \section{Spacetimes}
        My other interest lies in Lorentz geometry and spacetimes,
        the mathematical framework of general relativity. A Lorentz manifold
        is a pseudo-Riemannian manifold $(X,\,g)$ with signature $(n,\,1)$ for
        some positive integer $n$. That is, an $n+1$ dimensional smooth
        manifold together with a pseudo-Riemannian metric $g$%
        \footnote{
            Similar to a Riemannian metric, but the positive-definite
            criterion is relaxed to non-degeneracy.
        }
        such that for every point $p\in{X}$ there is a chart
        $(x_{1},\,\dots,\,x_{n},\,t)$ about $p$ such that:
        \begin{equation}
            g=-\textrm{d}t^{2}+\sum_{k=1}^{n}\textrm{d}x_{n}^{2}
        \end{equation}
        The null-vectors $v\in{T}_{x}X$, tangent vectors with
        $g_{x}(v,\,v)=0$, satisfy the equation of a cone. This
        \textit{light-cone} splits the tangent space into three parts: the
        cone, the exterior, and the interior, with the interior consisting
        of two connected components. A \textit{spacetime} is a Lorentz manifold
        $(X,\,g)$ with a smoothly varying choice of these connected
        components. This choice is called the \textit{future} direction at
        each point, and the result is a \textit{time-orientation} of the
        manifold. Spacetimes are the natural setting for the theory of
        general relativity, and the lines on these \textit{light cones} may be
        interpreted as \textit{rays of light}. Two points in a spacetime are
        \textit{causally related} if there is a piece-wise smooth curve
        $\gamma$ connecting them such that
        $g_{\gamma(t)}(\gamma'(t),\,\gamma'(t))\leq{0}$ for all $t$.
        Intuitively this means the speed of a particle traveling along $\gamma$
        never exceeds the speed of light.
        \par\hfill\par
        For appropriate spacetimes (globally hyperbolic) $X$, causality can be
        reduced to a question of \textit{linking}. The space of all light rays
        in a globally hyperbolic space corresponds naturally to the
        spherical cotangent bundle $ST^{*}M$ of a \textit{Cauchy surface} $M$,
        a particular codimension 1 submanifold of the spacetime, and the
        \textit{sky} of a point $p$ is given by the fiber of the canonical
        projection $\pi:ST^{*}M\rightarrow{M}$. Topologically these skies are
        spheres, and the question of causality can be reduced to
        \textit{Legendrian linking} of the skies of points. This is the
        Legendrian Low conjecture, proved by Chernov and Nemirovski in 2008
        \cite{ChernovNemirovski2010LowConj}.
        \par\hfill\par
        Knot theory has found its way into the study of spacetimes
        in other ways.
        The linking number, a famous invariant of links dating back to Gauss,
        has been generalized by the \textit{affine linking number} of the
        $(n-1)$ spheres in the spherical cotangent bundle of an $n$ dimensional
        manifold, and these spheres need not be zero homologous
        \cite{ChernovRudyak2003AffineLinkingGeneralTheory}.
        In \cite{ChernovRudyak2002AffineLinkingAndCausality}
        Chernov and Rudyak used this to prove a causal
        relation between points in a spacetime in regards to this invariant.
        Recently in 2024 Chernov and I used the notion of affine linking
        number to obtain estimates for the number of times an observer would
        see a star that is being gravitationally lensed
        \cite{MaguireChernovAffineLinkingNumber}.
    \section{Previous Work}
        The first project I worked on was DWEL, the
        Dual Wavelength Echidna LIDAR
        (\textbf{LI}ght \textbf{D}etectction \textbf{A}nd \textbf{R}anging),
        described in \cite{DWEL2012}.
        By using two lasers at 1064 nm and 1548 nm the instrument is able
        to distinguish scattering from leaves and tree trunks. This allows
        the mapping of forests for botanical studies. I, Glenn Howe, and
        Kuravi Hewawasum worked on this at the Center for Atmospheric
        Research (CAR) in Lowell, MA, in collaboration with UMass Boston
        and Boston University, and the device saw uses in California,
        Harvard forest, and Australia
        \cite{Li2016RadiometricCO}.
        \par\hfill\par
        The bulk of my work at CAR, renamed LoCSST
        (Lowell Center for Space Science and Technology), was on
        HiT\&MIS, a spectrometer used to study the aurora borealis
        (\cite{2011AGUFMSA13B1890C}, \cite{2014AGUFMSA13B4000H},
        \cite{2015AGUFMSA13B2369A}). The
        instrument used an echelle grating to achieve a high dispersion
        order allowing several wavelengths of light to be detected and
        differentiated between. Two copies of the instrument were built
        and deployed in parallel with Dartmouth's BARREL mission in
        Kiruna, Sweden. These efforts led to several talks and a
        publication \cite{AryalHewawasamMaguire2018DerivationHitAndMIS}.
        \par\hfill\par
        The final project I worked on in Lowell was the analysis of data
        from the SPINR sounding rocket mission. This is also where I began
        to work in programming and implementing algorithms. SPINR was a very
        interesting idea for a rocket, its detector only captured one
        spatial dimension (the second axis acted as a spectrograph). To
        create a two-dimensional picture the rocket spun about smoothly
        in the sky. The resulting \textit{sinograms} needed to be inverted
        to obtain an image of the night sky. This involved working with
        very large (26000-by-26000) matrices and delved into the theory of
        numerical linear algebra and sparse matrix representation. The
        final coding project was archived at MAST.
    \section{Future Work: Exoplanet Microlensing}
        My previous and current research has provided me technical skills in
        Fourier analysis, high-performance computing, and
        a particular love for planets beyond Earth.
        My theoretical background has also provided me with a knowledge of
        spacetimes, general relativity, and gravitational lensing.
        My new aim is to explore
        planets beyond Saturn and Uranus, and beyond our very own solar system.
        This past
        experience provides with the skills necessary to succeed in exoplanet
        microlensing and analyzing the data to detect new worlds.
    \bibliographystyle{plain}
    \bibliography{bib.bib}
\end{document}
